\documentclass[11pt]{article}
\usepackage[utf8]{inputenc}
\usepackage[T1]{fontenc}
\usepackage{amsmath}
\usepackage{amssymb}
\usepackage{multicol}
\usepackage{geometry}
\usepackage{tikz}
\usepackage{enumitem}
\usepackage{xcolor}
\usepackage{titlesec}

\geometry{a4paper, left=1cm, right=1cm, top=0.5cm, bottom=1.2cm}
\setlength{\columnseprule}{0.4pt}

\titleformat{\section}{\normalfont\Large\bfseries\color{blue}}{}{0em}{}
\titleformat{\subsection}{\normalfont\large\bfseries\color{red}}{}{0em}{}

\title{\textcolor{red}{Grandezas Diretamente Proporcionais}}
\author{Professor: Jefferson}
\date{}

\begin{document}

\maketitle
\vspace{-1cm}

\begin{multicols}{2}

\section*{Exercícios com Resolução Comentada}

\begin{enumerate}[leftmargin=*]

% QUESTÃO 1
\item Se 5 cadernos custam R\$ 40, quanto custarão 8 cadernos?

\textbf{Resolução:}

Como as grandezas são diretamente proporcionais:

\[
\frac{5}{40} = \frac{8}{x}
\]

\[
5x = 320
\Rightarrow x = 64
\]

\textbf{Resposta:} R\$ 64.

% QUESTÃO 2
\item Uma máquina produz 300 peças em 6 horas. Quantas peças produzirá em 10 horas?

\textbf{Resolução:}

\[
\frac{300}{6} = \frac{x}{10}
\]

\[
6x = 3000
\Rightarrow x = 500
\]

\textbf{Resposta:} 500 peças.

% QUESTÃO 3
\item Se 12 litros de combustível permitem percorrer 180 km, quantos quilômetros serão percorridos com 20 litros?

\textbf{Resolução:}

\[
\frac{12}{180} = \frac{20}{x}
\]

\[
12x = 3600
\Rightarrow x = 300
\]

\textbf{Resposta:} 300 km.

% QUESTÃO 4
\item Um pedreiro recebe R\$ 150 por dia. Quanto receberá em 18 dias?

\textbf{Resolução:}

\[
\frac{150}{1} = \frac{x}{18}
\]

\[
x = 150 \cdot 18 = 2700
\]

\textbf{Resposta:} R\$ 2.700.

% QUESTÃO 5
\item 4 torneiras enchem um reservatório em 2 horas. Mantendo a mesma vazão, quantas torneiras serão necessárias para encher 3 reservatórios no mesmo tempo?

\textbf{Resolução:}

Como aumenta o número de reservatórios, aumenta o número de torneiras.

\[
\frac{4}{1} = \frac{x}{3}
\]

\[
x = 12
\]

\textbf{Resposta:} 12 torneiras.

% QUESTÃO 6
\item 7 operários constroem um muro em 5 dias. Quantos operários são necessários para construir 3 muros no mesmo período?

\textbf{Resolução:}

\[
\frac{7}{1} = \frac{x}{3}
\]

\[
x = 21
\]

\textbf{Resposta:} 21 operários.

% QUESTÃO 7
\item Uma receita utiliza 250 g de farinha para fazer um bolo. Quantos gramas serão necessários para 6 bolos?

\textbf{Resolução:}

\[
\frac{250}{1} = \frac{x}{6}
\]

\[
x = 1500
\]

\textbf{Resposta:} 1500 g.

% QUESTÃO 8
\item 3 kg de carne servem 6 pessoas. Quantos kg serão necessários para 15 pessoas?

\textbf{Resolução:}

\[
\frac{3}{6} = \frac{x}{15}
\]

\[
6x = 45
\Rightarrow x = 7{,}5
\]

\textbf{Resposta:} 7,5 kg.

% QUESTÃO 9
\item Se 8 metros de tecido custam R\$ 200, quanto custarão 15 metros?

\textbf{Resolução:}

\[
\frac{8}{200} = \frac{15}{x}
\]

\[
8x = 3000
\Rightarrow x = 375
\]

\textbf{Resposta:} R\$ 375.

% QUESTÃO 10
\item Uma gráfica imprime 2.400 panfletos em 3 horas. Quantos panfletos imprimirá em 8 horas?

\textbf{Resolução:}

\[
\frac{2400}{3} = \frac{x}{8}
\]

\[
3x = 19200
\Rightarrow x = 6400
\]

\textbf{Resposta:} 6.400 panfletos.

\end{enumerate}

\end{multicols}

\end{document}
